% Generated by tools/edn_latex.py from reports/edn.md.
% Do not edit: change reports/edn.md and run `make edn`.
\ednentry{
  id = {EDN-26},
  title = {NFR-11's window is 1,000 requests, and \texorpdfstring{\texttt{model}}{model} is excluded from the drift comparison},
  note = {\textbf{Amended by EDN-29} (2026-09-29): the drift job runs at a significance level of 0.005, not 0.05.
The configuration below was measured on a scope that EDN-22 has since narrowed, and on one split; re-measured, its
control clause does not hold at 0.05.},
  date = {2026-09-29},
  milestone = {M6: Monitoring},
  topic = {Monitoring, Requirements},
  participants = {@lukas2510},
  decision = {Amends EDN-14. The drift job compares windows of \textbf{1,000} requests, not 100, against the 10,000-listing reference, and leaves \textbf{\texttt{model}} out of the comparison. The other clauses of NFR-11 are unchanged: the \texttt{ES} replay must be flagged, at least three features other than \texttt{country\_\allowbreak{}code} must be named, and at most 1 of 20 i.i.d. control windows may be flagged.},
  rationale = {\emph{Alternatives considered:}
\begin{itemize}[leftmargin=*, topsep=2pt]
\item \textbf{Option A (chosen): window 1,000 and \texttt{model} excluded.}
\newline Pros: the only configuration measured to satisfy all three clauses as written. Both changes are independently justified: three drifted features do not reliably co-occur in a 100-row window, and a 349-level chi-square against such a window is under-powered rather than merely noisy.
\newline Cons: the M6 drift demonstration needs 1,000 replayed requests instead of 100, so the monitoring section of the report shows a slower signal.
\item \textbf{Option B: window 1,000, keep \texttt{model}, relax the control clause to at most 2 in 20.}
\newline Pros: no feature is removed, so the monitoring watches everything it sees.
\newline Cons: weakens a promise instead of fixing its cause, and keeps a test whose cell counts do not support the chi-square.
\item \textbf{Option C: keep window 100 and relax both clauses (at least one changed feature, at most 3 false alarms in 20).}
\newline Pros: keeps the fast 100-request demonstration.
\newline Cons: two promises weakened at once; 3 false alarms in 20 windows is hard to defend for a monitoring requirement.
\item \textbf{Option D: keep NFR-11 as EDN-14 wrote it.}
\newline Cons: measured to fail. At window 100 with all features the control flags 5.4 of 20 windows, and the three-feature clause holds in only 81.5\,\% of trials.
\end{itemize}
\par
\emph{Why:} NFR-11 as written was never measured in one configuration: EDN-14 took its \texttt{ES} half from \texttt{nfr11\_\allowbreak{}check.\allowbreak{}py} (window 100, full reference) and its control half from \texttt{nfr11\_\allowbreak{}diag.\allowbreak{}py} (window 1,000, reference capped at 10,000), and the latter's output was never committed. Measuring both halves together (\texttt{nfr11\_\allowbreak{}model\_\allowbreak{}excluded.\allowbreak{}py}, 200 trials per cell, 10,000-listing reference, i.i.d. control) gives:
\par
{\small\renewcommand{\arraystretch}{1.15}%
\begin{tabularx}{\linewidth}{@{}>{\raggedright\arraybackslash}X>{\raggedright\arraybackslash}X>{\raggedright\arraybackslash}X>{\raggedright\arraybackslash}X@{}}
\toprule
Configuration & \texttt{ES} flagged & three features besides \texttt{country\_\allowbreak{}code} & control windows flagged of 20 \\
\midrule
all features, window 100 & 1.000 & 81.5\,\% of trials, minimum 1 & 5.4 \\
\texttt{model} excluded, window 100 & 1.000 & 79.5\,\% of trials, minimum 1 & 2.2 \\
all features, window 1,000 & 1.000 & 100\,\%, minimum 13 & 1.4 \\
\textbf{\texttt{model} excluded, window 1,000} & \textbf{1.000} & \textbf{100\,\%, minimum 12} & \textbf{0.7} \\
\bottomrule
\end{tabularx}}
\par
Neither change suffices alone. Excluding \texttt{model} roughly halves the false alarms at both window sizes, because at window 100 it alone accounts for 0.190 of the 0.270 flag rate, but it does not reach 1 in 20 and does not help the three-feature clause at all. Enlarging the window fixes the three-feature clause and most of the false alarms, but with \texttt{model} kept the control still flags 1.4 of 20. Only the combination satisfies every clause, so NFR-11 now states a configuration that has been demonstrated rather than assembled from two runs.},
  ai = {Information seeking, Alternative generation, Alternative assessment, Recommendation, Solution generation},
  response = {Accepted with modifications},
  assessment = {A reviewing agent reported that the control clause fails at window 100 and proposed stating window 1,000. A second agent argued that the window was the wrong lever, because the false alarms come from \texttt{model}, and recommended excluding it while keeping window 100; Lukas chose that option. The measurement then refuted it: the control improves from 5.4 to 2.2 flagged windows of 20 but stays above the promised 1, and the three-feature clause fails in one run in five at window 100 either way. Two claims in that recommendation were also wrong and are recorded rather than quietly dropped: \texttt{model} was described as having roughly 10,000 levels, where it has 349 in the reference (84,171 is \texttt{model\_\allowbreak{}version}, a different column), and the reviewing agent's control figure at window 1,000 was 0.6 of 20 against the 1.4 measured here, a difference that decides whether the clause holds with \texttt{model} kept. The argument that a 349-level chi-square against a 100-row window is under-powered survived the measurement; its effect was simply smaller than claimed. Lukas then chose the combination, which is the only configuration that was demonstrated to work.},
  aievidence = {Claude Code session on 2026-09-29: a peer agent's review of PR \#19 reported the control failure and proposed window 1,000; this session verified the finding by inspection, argued for excluding \texttt{model} instead, and Lukas chose that. After the measurement refuted it, this session reported the four measured configurations and recommended the combination, which Lukas chose.},
  evidence = {\href{https://github.com/mlops-2627q1-mds-upc/MLOPS_Recommenditos/blob/main/docs/docs/requirements.md}{requirements} NFR-11; \href{https://github.com/mlops-2627q1-mds-upc/MLOPS_Recommenditos/blob/main/docs/docs/specification.md}{specification} NFR-11; \href{https://github.com/mlops-2627q1-mds-upc/MLOPS_Recommenditos/blob/main/reports/analysis}{reports/analysis/} (\texttt{nfr11\_\allowbreak{}model\_\allowbreak{}excluded.\allowbreak{}py}, \texttt{nfr11\_\allowbreak{}model\_\allowbreak{}excluded\_\allowbreak{}results.\allowbreak{}json}, run of 2026-09-29); \hyperref[EDN-14]{EDN-14}; \href{https://github.com/mlops-2627q1-mds-upc/MLOPS_Recommenditos/pull/19}{PR \#19}.},
}
